\subsection{混合边界离散 Laplacian 的半奇整数量子化与四联谱刚性}\label{subsec:conclusion-halfodd-laplacian-bulkboundary-bessel-oddprojector-cosh}
沿用混合边界离散 Laplacian
\[
L_k:=K_k^{-1},
\qquad
\mu_p
=4\sin^2\!\frac{(2p-1)\pi}{4k+2}
=2\Bigl(1-\cos\frac{(2p-1)\pi}{2k+1}\Bigr),
\qquad p=1,\dots,k,
\]
以及经验谱测度与端点谱测度
\[
\nu_k:=\frac1k\sum_{p=1}^k\delta_{\mu_p},
\qquad
\rho_k:=\sum_{p=1}^k w_{p,k}\,\delta_{\mu_p},
\qquad
w_{p,k}=\frac{4}{2k+1}\sin^2\theta_p,
\qquad
\theta_p:=\frac{(2p-1)\pi}{2k+1}.
\]
此前结果已经分别给出：\(\nu_k\) 的反正弦极限、\(\rho_k\) 的 \(\mathrm{Beta}(\tfrac32,\tfrac32)\) 极限、离散 Weyl 计数的显式取整式、热核的 Bessel 闭式、规范化谱 \(\zeta\) 的 odd-projector 极限，以及 ND 行列式对 \(\cosh(\sqrt{s})\) 的局部一致收敛。将这些接口合并之后，可把半奇整数量子化机制压缩为一条同时支配 bulk 谱、boundary 谱、热核、谱 \(\zeta\) 与 entire 行列式的统一闭链。

\begin{theorem}[bulk--boundary 极限谱测度的二次倾斜律]\label{thm:conclusion-Lk-bulk-boundary-quadratic-tilt}
当 \(k\to\infty\) 时，
\[
\nu_k\Rightarrow \nu_\infty,
\qquad
\dd\nu_\infty(\mu)=\frac{1}{\pi\sqrt{\mu(4-\mu)}}\,\dd\mu,
\qquad 0<\mu<4,
\]
并且
\[
\rho_k\Rightarrow \rho_\infty,
\qquad
\dd\rho_\infty(\mu)=\frac{\sqrt{\mu(4-\mu)}}{2\pi}\,\dd\mu,
\qquad 0<\mu<4.
\]
更强地，二者满足精确倾斜关系
\[
\dd\rho_\infty(\mu)=\frac{\mu(4-\mu)}{2}\,\dd\nu_\infty(\mu).
\]
亦即，端点谱测度正是 bulk 反正弦律经二次倾斜 \(\mu(4-\mu)/2\) 所得到的 boundary 极限。
\end{theorem}

\begin{proof}
前两式分别由定理 \ref{thm:pom-Lk-arcsine-law} 与定理 \ref{thm:pom-Lk-boundary-spectral-measure-beta32} 给出。
最后一式正是推论 \ref{cor:pom-Lk-boundary-bulk-radon-nikodym}。
\end{proof}

\begin{corollary}[中心二项式与 Catalan 的双矩律]\label{cor:conclusion-Lk-central-binomial-catalan-moments}
\leanverified{paper\_conclusion\_lk\_central\_binomial\_catalan\_moments}
对任意整数 \(r\ge 0\)，有
\[
\int_0^4 \mu^r\,\dd\nu_\infty(\mu)=\binom{2r}{r},
\qquad
\int_0^4 \mu^r\,\dd\rho_\infty(\mu)=\frac{1}{r+2}\binom{2r+2}{r+1}.
\]
因此，bulk 极限矩由中心二项式系数给出，而 boundary 极限矩恰为偏移一位的 Catalan 数。
\end{corollary}

\begin{proof}
对 \(\nu_\infty\) 作代换 \(\mu=4x\) 并使用 Beta 函数恒等式，得到
\[
\int_0^4 \mu^r\,\dd\nu_\infty(\mu)
=
4^r\frac{B(r+\tfrac12,\tfrac12)}{B(\tfrac12,\tfrac12)}
=\binom{2r}{r}.
\]
对 \(\rho_\infty\) 同理有
\[
\int_0^4 \mu^r\,\dd\rho_\infty(\mu)
=
4^r\frac{B(r+\tfrac32,\tfrac32)}{B(\tfrac32,\tfrac32)}
=\frac{1}{r+2}\binom{2r+2}{r+1}.
\]
第二式也可直接视为推论 \ref{cor:pom-Lk-boundary-catalan-moments-gf} 的矩封闭式。
\end{proof}

\begin{corollary}[端点加权把硬边奇性抬升一阶]\label{cor:conclusion-Lk-hard-edge-lift}
当 \(\varepsilon\downarrow 0\) 时，
\[
\nu_\infty([0,\varepsilon])
=
\frac{1}{\pi}\sqrt{\varepsilon}+O(\varepsilon^{3/2}),
\qquad
\rho_\infty([0,\varepsilon])
=
\frac{2}{3\pi}\varepsilon^{3/2}+O(\varepsilon^{5/2}).
\]
在上端点 \(4\) 处也有对称展开
\[
\nu_\infty([4-\varepsilon,4])
=
\frac{1}{\pi}\sqrt{\varepsilon}+O(\varepsilon^{3/2}),
\qquad
\rho_\infty([4-\varepsilon,4])
=
\frac{2}{3\pi}\varepsilon^{3/2}+O(\varepsilon^{5/2}).
\]
因此，端点权把 bulk 反正弦律的平方根硬边奇性，严格提升为三次半幂边界层。
\end{corollary}

\begin{proof}
由定理 \ref{thm:conclusion-Lk-bulk-boundary-quadratic-tilt} 的密度表达，
\[
\dd\nu_\infty(\mu)\sim \frac{1}{2\pi}\mu^{-1/2}\,\dd\mu,
\qquad
\dd\rho_\infty(\mu)\sim \frac{1}{\pi}\mu^{1/2}\,\dd\mu
\qquad(\mu\downarrow 0).
\]
积分即得
\[
\nu_\infty([0,\varepsilon])=\frac{1}{\pi}\sqrt{\varepsilon}+O(\varepsilon^{3/2}),
\qquad
\rho_\infty([0,\varepsilon])=\frac{2}{3\pi}\varepsilon^{3/2}+O(\varepsilon^{5/2}).
\]
上端点 \(4\) 的情形由变换 \(\mu\mapsto 4-\mu\) 的对称性完全相同。
\end{proof}

\begin{theorem}[离散 Weyl 律导出的 sharp \texorpdfstring{$1/k$}{1/k} Kolmogorov 界]\label{thm:conclusion-Lk-kolmogorov-sharp-1overk}
\leanverified{paper\_conclusion\_lk\_kolmogorov\_sharp\_1overk}
记
\[
F_k(\mu):=\nu_k((-\infty,\mu])=\frac{N_k(\mu)}{k},
\qquad
F_\infty(\mu):=\nu_\infty((-\infty,\mu])=\frac{\theta(\mu)}{\pi},
\]
其中
\[
\theta(\mu):=\arccos\!\Bigl(1-\frac{\mu}{2}\Bigr).
\]
则对所有 \(\mu\in[0,4]\)，有统一界
\[
\bigl|F_k(\mu)-F_\infty(\mu)\bigr|\le \frac{1}{k}.
\]
因此，\(\nu_k\) 对其连续极限的 Kolmogorov 距离以 sharp 的 \(O(1/k)\) 速度衰减。
\end{theorem}

\begin{proof}
由推论 \ref{cor:pom-Lk-discrete-weyl}，
\[
N_k(\mu)=\left\lfloor \frac{2k+1}{2\pi}\theta(\mu)+\frac12\right\rfloor.
\]
记
\[
A(\mu):=\frac{2k+1}{2\pi}\theta(\mu).
\]
则 \(|N_k(\mu)-A(\mu)|\le \frac12\)，故
\[
\left|\frac{N_k(\mu)}{k}-\frac{\theta(\mu)}{\pi}\right|
\le
\frac{1}{2k}
+
\left|\frac{A(\mu)}{k}-\frac{\theta(\mu)}{\pi}\right|
=
\frac{1}{2k}+\frac{\theta(\mu)}{2\pi k}
\le \frac{1}{k},
\]
因为 \(0\le \theta(\mu)\le \pi\)。
\end{proof}

\begin{corollary}[固定量子编号的硬边极限]\label{cor:conclusion-Lk-fixed-index-hard-edge}
对每个固定 \(p\ge 1\)，成立
\[
(2k+1)^2\mu_p\longrightarrow (2p-1)^2\pi^2.
\]
更精确地，
\[
\mu_p=
\frac{(2p-1)^2\pi^2}{(2k+1)^2}+O(k^{-4})
\qquad
(k\to\infty,\ p\ \text{固定}).
\]
\end{corollary}

\begin{proof}
由
\[
\mu_p=4\sin^2\!\Bigl(\frac{(2p-1)\pi}{4k+2}\Bigr)
\]
及 \(\sin x=x+O(x^3)\)，得
\[
\mu_p=
4\left(\frac{(2p-1)\pi}{4k+2}\right)^2+O(k^{-4})
=
\frac{(2p-1)^2\pi^2}{(2k+1)^2}+O(k^{-4}).
\qedhere
\]
\end{proof}

\begin{theorem}[bulk--boundary 热核的 Bessel 对偶]\label{thm:conclusion-Lk-bulk-boundary-bessel-duality}
\leanverified{paper\_conclusion\_lk\_bulk\_boundary\_bessel\_duality}
对任意 \(t>0\)，有
\[
\int_0^4 e^{-t\mu}\,\dd\nu_\infty(\mu)=e^{-2t}I_0(2t),
\]
以及
\[
\int_0^4 e^{-t\mu}\,\dd\rho_\infty(\mu)
=
e^{-2t}\bigl(I_0(2t)-I_2(2t)\bigr)
=
e^{-2t}\frac{I_1(2t)}{t}.
\]
因此，boundary 热核正是同一 Bessel 家族中的下一阶闭式像。
\end{theorem}

\begin{proof}
第一式是推论 \ref{cor:pom-Lk-heat-kernel-bessel} 的连续极限表达。
对第二式，由定理 \ref{thm:pom-Lk-boundary-spectral-measure-beta32} 的角变量表示，
\[
\int_0^4 e^{-t\mu}\,\dd\rho_\infty(\mu)
=
\frac{2}{\pi}\int_0^\pi \sin^2\theta\,e^{-2t(1-\cos\theta)}\,\dd\theta.
\]
再利用
\[
\sin^2\theta=\frac{1-\cos 2\theta}{2},
\qquad
\frac{1}{\pi}\int_0^\pi e^{2t\cos\theta}\cos(n\theta)\,\dd\theta=I_n(2t),
\]
得到
\[
\int_0^4 e^{-t\mu}\,\dd\rho_\infty(\mu)
=
e^{-2t}\bigl(I_0(2t)-I_2(2t)\bigr).
\]
最后应用递推恒等式 \(I_0(x)-I_2(x)=2I_1(x)/x\) 并取 \(x=2t\)，即得
\[
e^{-2t}\bigl(I_0(2t)-I_2(2t)\bigr)=e^{-2t}\frac{I_1(2t)}{t}.
\]
\end{proof}

\begin{corollary}[边界热核是 bulk 热核的二阶差分像]\label{cor:conclusion-Lk-boundary-heat-as-bulk-differential}
设
\[
F(t):=\int_0^4 e^{-t\mu}\,\dd\nu_\infty(\mu)=e^{-2t}I_0(2t).
\]
则对任意 \(t>0\)，有
\[
\int_0^4 e^{-t\mu}\,\dd\rho_\infty(\mu)
=
-2F'(t)-\frac12F''(t).
\]
\end{corollary}

\begin{proof}
由定理 \ref{thm:conclusion-Lk-bulk-boundary-quadratic-tilt}，
\[
\int_0^4 e^{-t\mu}\,\dd\rho_\infty(\mu)
=
\frac12\int_0^4 \mu(4-\mu)e^{-t\mu}\,\dd\nu_\infty(\mu).
\]
另一方面，
\[
F'(t)=-\int_0^4 \mu e^{-t\mu}\,\dd\nu_\infty(\mu),
\qquad
F''(t)=\int_0^4 \mu^2 e^{-t\mu}\,\dd\nu_\infty(\mu),
\]
故
\[
\frac12\int_0^4 \mu(4-\mu)e^{-t\mu}\,\dd\nu_\infty(\mu)
=
\frac12\int_0^4 (4\mu-\mu^2)e^{-t\mu}\,\dd\nu_\infty(\mu)
=
-2F'(t)-\frac12F''(t).
\]
\end{proof}

\begin{theorem}[规范化谱 \texorpdfstring{\(\zeta\)}{zeta} 的 odd-projector 极限]\label{thm:conclusion-Lk-normalized-zeta-odd-projector}
定义
\[
\widetilde\zeta_k(s):=
\left(\frac{\pi}{2k+1}\right)^{2s}\zeta_k(s),
\qquad
\zeta_k(s):=\sum_{p=1}^k \mu_p^{-s}.
\]
则对任意紧集 \(K\subset\{\,\Re(s)>\tfrac12\,\}\)，一致成立
\[
\widetilde\zeta_k(s)
=
\bigl(1-2^{-2s}\bigr)\zeta(2s)+O_K\!\bigl(k^{1-2\inf_{s\in K}\Re(s)}\bigr).
\]
特别地，
\[
\widetilde\zeta_k(s)\longrightarrow \bigl(1-2^{-2s}\bigr)\zeta(2s)
\qquad
\text{于 }K\text{ 上一致}.
\]
因此，半奇整数量子化在规范化谱 \(\zeta\) 中强制产生奇模态 Dirichlet 投影因子 \((1-2^{-2s})\)。
\end{theorem}

\begin{proof}
由命题 \ref{prop:pom-Lk-spectral-zeta-dirichlet}，
\[
\zeta_k(s)
=
\left(\frac{2k+1}{\pi}\right)^{2s}\bigl(1-2^{-2s}\bigr)\zeta(2s)+O(k)
\]
在 \(\Re(s)>\tfrac12\) 上成立。
两边同乘 \(\bigl(\pi/(2k+1)\bigr)^{2s}\) 得
\[
\widetilde\zeta_k(s)
=
\bigl(1-2^{-2s}\bigr)\zeta(2s)
+
O(k)\left(\frac{\pi}{2k+1}\right)^{2s}.
\]
若 \(K\subset\{\Re(s)>\tfrac12\}\) 为紧集，则
\[
\sup_{s\in K}\left|O(k)\left(\frac{\pi}{2k+1}\right)^{2s}\right|
=
O_K\!\bigl(k^{1-2\inf_{s\in K}\Re(s)}\bigr)\to 0,
\]
从而得到局部一致收敛。
\end{proof}

\begin{corollary}[规范化 hard-edge 谱 \texorpdfstring{\(\zeta\)}{zeta} 在 \texorpdfstring{$s=\tfrac12$}{s=1/2} 的留数]\label{cor:conclusion-Lk-normalized-zeta-residue-quarter}
函数
\[
\widetilde\zeta_\infty(s):=\bigl(1-2^{-2s}\bigr)\zeta(2s)
\]
在 \(s=\tfrac12\) 处具有简单极点，并满足
\[
\operatorname*{Res}_{s=1/2}\widetilde\zeta_\infty(s)=\frac14.
\]
\end{corollary}

\begin{proof}
\(\zeta(u)\) 在 \(u=1\) 处的留数为 \(1\)，故作变量代换 \(u=2s\) 得
\[
\operatorname*{Res}_{s=1/2}\zeta(2s)=\frac12.
\]
又
\[
\left(1-2^{-2s}\right)\Big|_{s=1/2}=\frac12,
\]
故所求留数为 \(\frac12\cdot\frac12=\frac14\)。
\end{proof}

\begin{theorem}[ND 行列式极限的零点逐一追踪]\label{thm:conclusion-Lk-nd-determinant-zero-tracking}
定义
\[
D_k(s):=\det\!\left(L_k+\frac{4s}{(2k+1)^2}I\right).
\]
则 \(D_k(s)\to \cosh(\sqrt{s})\) 在 \(s\)-平面的每个紧集上一致成立。
并且 \(D_k\) 的全部零点恰为
\[
s_{k,p}:=-\frac{(2k+1)^2}{4}\mu_p,
\qquad 1\le p\le k,
\]
且对每个固定 \(p\ge 1\)，有
\[
s_{k,p}\longrightarrow -\left(\frac{(2p-1)\pi}{2}\right)^2.
\]
后者正是 \(\cosh(\sqrt{s})\) 的全部零点。
\end{theorem}

\begin{proof}
局部一致收敛 \(D_k\to \cosh(\sqrt{s})\) 正是推论 \ref{cor:pom-Lk-continuum-ND-det-cosh-locally-uniform} 的内容。
另一方面，
\[
D_k(s)=\prod_{p=1}^k\left(\mu_p+\frac{4s}{(2k+1)^2}\right),
\]
故其零点正是所给 \(s_{k,p}\)。
再由推论 \ref{cor:conclusion-Lk-fixed-index-hard-edge}，
\[
s_{k,p}
=
-\frac{(2k+1)^2}{4}\mu_p
\longrightarrow
-\frac14(2p-1)^2\pi^2
=
-\left(\frac{(2p-1)\pi}{2}\right)^2.
\]
而 \(\cosh z=0\) 当且仅当 \(z=i(2q-1)\pi/2\)，故 \(\cosh(\sqrt{s})\) 的零点恰为上述极限点。
\end{proof}

\begin{corollary}[第一零点的单调逼近律]\label{cor:conclusion-Lk-first-zero-monotone-approach}
定义
\[
\alpha_k:=\frac{(2k+1)^2}{4}\mu_1
=(2k+1)^2\sin^2\!\frac{\pi}{4k+2}.
\]
则 \((\alpha_k)_{k\ge 1}\) 严格递增，且
\[
\alpha_k\uparrow \frac{\pi^2}{4}.
\]
等价地，
\[
\frac{\lambda_{\max}(K_k)}{(2k+1)^2}\downarrow \frac{1}{\pi^2}.
\]
\end{corollary}

\begin{proof}
记
\[
x_k:=\frac{\pi}{4k+2},
\qquad
\alpha_k=\frac{\pi^2}{4}\left(\frac{\sin x_k}{x_k}\right)^2.
\]
由于 \(x_k\downarrow 0\)，而函数 \(x\mapsto \sin x/x\) 在 \((0,\pi)\) 上严格递减，故 \(\alpha_k\) 严格递增。
再由 \(\sin x/x\to 1\) 得 \(\alpha_k\to \pi^2/4\)。
最后，\(\lambda_{\max}(K_k)=\mu_1^{-1}\)，故
\[
\frac{\lambda_{\max}(K_k)}{(2k+1)^2}
=
\frac{1}{(2k+1)^2\mu_1}
=
\frac{1}{4\alpha_k}
\downarrow \frac{1}{\pi^2}.
\]
\end{proof}

\begin{corollary}[归一化迹矩的中心二项式极限]\label{cor:conclusion-Lk-normalized-trace-moments-central-binomial}
对任意整数 \(r\ge 0\)，有
\[
\lim_{k\to\infty}\frac{1}{k}\operatorname{tr}(L_k^r)
=
\binom{2r}{r}.
\]
\end{corollary}

\begin{proof}
由经验谱测度的定义，
\[
\frac{1}{k}\operatorname{tr}(L_k^r)=\int_0^4 \mu^r\,\dd\nu_k(\mu).
\]
再由定理 \ref{thm:pom-Lk-arcsine-law} 的弱收敛与推论 \ref{cor:conclusion-Lk-central-binomial-catalan-moments} 的 bulk 矩公式，得到
\[
\lim_{k\to\infty}\frac{1}{k}\operatorname{tr}(L_k^r)
=
\int_0^4 \mu^r\,\dd\nu_\infty(\mu)
=
\binom{2r}{r}.
\]
\end{proof}

\begin{theorem}[半奇整数混合边界谱的四联刚性]\label{thm:conclusion-halfodd-mixed-boundary-spectrum-fourfold-rigidity}
设
\[
\mu_p^{(k)}=4\sin^2\!\frac{(2p-1)\pi}{4k+2},
\qquad
L_k=K_k^{-1}.
\]
则由同一半奇整数量子化骨架同时强制出以下四类连续对象：
\begin{enumerate}
  \item bulk 与 boundary 的极限谱测度分别为
  \[
  \dd\nu_\infty(\mu)=\frac{1}{\pi\sqrt{\mu(4-\mu)}}\,\dd\mu,
  \qquad
  \dd\rho_\infty(\mu)=\frac{\sqrt{\mu(4-\mu)}}{2\pi}\,\dd\mu,
  \]
  且满足精确倾斜律
  \[
  \dd\rho_\infty(\mu)=\frac{\mu(4-\mu)}{2}\,\dd\nu_\infty(\mu);
  \]
  \item bulk 与 boundary 热核密度分别为
  \[
  F(t)=e^{-2t}I_0(2t),
  \qquad
  G(t)=e^{-2t}\bigl(I_0(2t)-I_2(2t)\bigr)=e^{-2t}\frac{I_1(2t)}{t};
  \]
  \item 规范化谱 \(\zeta\) 满足
  \[
  \zeta_k(s)\sim \left(\frac{2k+1}{\pi}\right)^{2s}\!\bigl(1-2^{-2s}\bigr)\zeta(2s)
  \qquad
  (\Re s>\tfrac12);
  \]
  \item ND 行列式满足局部一致收敛
  \[
  D_k(s)\longrightarrow \cosh(\sqrt{s}),
  \]
  且其零点极限恰追踪到
  \[
  -\left(\frac{(2p-1)\pi}{2}\right)^2.
  \]
\end{enumerate}
因此，混合边界谱并不分裂为互不相干的“测度 / 热核 / \(\zeta\) / 行列式”四层；它们只是同一半奇整数骨架在 Stieltjes、Laplace、Mellin 与 entire 行列式语义下的四重投影。
\end{theorem}

\begin{proof}
第 \(1\) 条由定理 \ref{thm:conclusion-Lk-bulk-boundary-quadratic-tilt} 给出。
第 \(2\) 条由定理 \ref{thm:conclusion-Lk-bulk-boundary-bessel-duality} 给出。
第 \(3\) 条是定理 \ref{thm:conclusion-Lk-normalized-zeta-odd-projector} 的改写。
第 \(4\) 条则正是定理 \ref{thm:conclusion-Lk-nd-determinant-zero-tracking}。
这些对象都只依赖同一族半奇整数量子化谱点 \(\mu_p^{(k)}\)，故四者并非独立输入，而是同一离散谱骨架在不同解析变换下的同步外显。证毕。
\end{proof}

\begin{conjecture}[半奇整数量子化的四联刚性]\label{conj:conclusion-halfodd-quantization-fourfold-rigidity}
设 \((\mathcal L_k)\) 为一族有限窗自伴离散模型，并满足以下两项结构条件：
\begin{enumerate}
\item 谱点具有半奇整数量子化形式
\[
\mu_{k,p}=2\bigl(1-\cos\theta_{k,p}\bigr),
\qquad
\theta_{k,p}=\frac{(2p-1)\pi}{2k+\beta};
\]
\item 存在与端点向量相容的加权谱测度极限。
\end{enumerate}
则该族模型应同时强制四类连续对象的联锁出现：
\[
\text{(i) bulk 反正弦律},\qquad
\text{(ii) boundary }\mathrm{Beta}\!\left(\tfrac32,\tfrac32\right)\text{ 律},
\]
\[
\text{(iii) 规范化谱 \(\zeta\) 中的 odd-projector 因子 }(1-2^{-2s}),
\]
\[
\text{(iv) ND 型行列式极限 }\cosh(\sqrt{s})\ \text{或其同类 entire 模型}.
\]
换言之，半奇整数量子化不应只是一条特征值参数化公式，而应是一种同时锁定 bulk 测度、boundary 测度、奇模态 Dirichlet 投影与 entire 行列式极限的统一刚性机制。
\end{conjecture}

\begin{remark}[统一读法]\label{rem:conclusion-halfodd-laplacian-bulkboundary-bessel-oddprojector-cosh}
定理 \ref{thm:conclusion-Lk-bulk-boundary-quadratic-tilt} 到定理 \ref{thm:conclusion-Lk-bulk-boundary-bessel-duality} 表明，bulk 与 boundary 两类谱测度并非彼此平行，而是通过二次倾斜因子 \(\mu(4-\mu)/2\) 精确耦合，进而把中心二项式矩、Catalan 矩与 Bessel 热核闭式纳入同一链条。
定理 \ref{thm:conclusion-Lk-normalized-zeta-odd-projector}、推论 \ref{cor:conclusion-Lk-normalized-zeta-residue-quarter} 与定理 \ref{thm:conclusion-Lk-nd-determinant-zero-tracking} 则表明，同一半奇整数量子化机制还会同时支配规范化谱 \(\zeta\) 的 odd-projector 极限与 ND 行列式的 \(\cosh(\sqrt{s})\) entire 极限。
因此，这里真正稳定的对象不是单个极限公式，而是一组被同一量子化规律共同锁定的四联谱刚性。
\end{remark}

\endinput
